% !TEX root = main.tex
\section{전이 검증과 참조 회계 관계}

\subsection{규칙 선택과 초기 설정}
이벤트의 검증 경로는 Table 2와 같이 구분한다. 일반 일곱 이
벤트는 배포 시 고정한 관계와 증명 검증기(verifier)를 사용한
다. 가공은 등록된 정책과 적용 범위별 사용 권한을 검사하고,
발급은 등록된 발급 정책을 검사하되 적용 범위별 사용 권한을
요구하지 않는다. 각 경로는 최종 영지식증명 하나를 사용한다.
공통 암호·생애주기 관계를 \(R_{ckore}\)라 하면 전체 회로 관계는

\begin{equation}
    R_k=R_k^{\mathrm{core}}\land R_k^{\mathrm{acct}}\tag{12}
\end{equation}
로 표현할 수 있다. 고정 경로의 회계 관계는 배포 설정으로 정
하고, 등록 경로에서는 선택한 정책 \(p\)의 관계로 정한다.

공통 관계에는 개방값과 커미트먼트, 포함 관계, 소유권 비
밀, 소비 식별자, 정확한 입출력 형상, 중복 입력·출력 금지, 출
력 및 감사 암호문의 결합과 정수 범위 검사가 포함된다. 현재
상태와 소비 여부는 공개 소비 식별자를 사용하여 검증 계약이

% Source PDF page 7, column 2
검사한다. 등록 경로의 \(p\)와 가공의 적용 범위는 증명에 결합하
고, 고정 경로의 이벤트 종류는 호출 경로와 고정 검증기 선택으
로도 결합할 수 있다.

배포자는 암호·직렬화·원장 문맥·역할·위원회 설정과 고정
검증기를 설치한다. 초기 회계 증표·이전 증표 트리는 비어 있
고 초기 루트는 수락된 루트이며, 소비·생성 기록·제품 주장·감
사 기록도 빈 상태에서 시작한다. 정책 관리자는 이벤트 종류,
정확한 입출력 수, 정책 이름과 버전으로 식별한 정책에 대해 검
증 관계 명세와 회로·검증키의 대응을 공개한다. 가공의 각 지
원 \((m,n)\)에는 별도 정책과 검증키를 사용한다. 정책 등록부의
검증기 메타데이터는 변경할 수 없고, 정책 활성값은 참에서 거
짓으로만 바꿀 수 있다. 비활성화는 이후 제출에 적용하며 과거
수락 이벤트나 제품 주장을 소급하여 무효화하지 않는다.

가공의 적용 범위 식별자 scope는 회사나 공정 배합을 직접
나타내지 않는 불투명 식별자다. 제출자 또는 자격증명의 적용
범위 권한과 해당 \((scope,p)\) 사용 권한을 함께 검사하며, 소유
권 지식만으로 모든 정책의 사용 권한을 주지 않는다. 고정 이벤
트의 감사 기록에는 정책 표식 0을, 가공과 발급에는 실제 \(p\)를
기록한다. 등록 정책 식별자는 이 표식 0과 구별한다.

\subsection{증명과 원자적 수락}
증명자는 실제 입력 개방값과 포함 증명 경로를 준비하고, 입력
소비 식별자와 출력 커미트먼트·미래 식별자를 계산한다. 실제
부모 참조와 출력 미래 식별자를 암호화한 뒤, 이벤트별 전체 관
계에 대한 증명을 생성한다. 공개 입력에는 필요한 루트, 소비
식별자, 순서 있는 출력과 감사 암호문, 해당 경로의 정책·적용
범위, 기한 등 이벤트별 메타데이터가 포함된다.

검증 계약은 선택된 검증기와 정확한 형상, 해당 정책·권한,
수락된 포함 관계 트리의 루트, 현재 미소비·Active 상태, 실행
시점의 기한과 모든 출력 키의 신규성을 검사한다. 동일 이벤트
안에서도 실제 입력 참조와 타입별 소비 식별자, 출력 객체 조회
키는 각각 서로 달라야 한다. 이미 등록된 제품 주장이나 기존
출력 키를 다시 생성할 수 없다.

모든 조건과 증명 검증이 통과하면 하나의 트랜잭션에서
aid를 배정하고, 소비 식별자에 소비 이벤트를 기록하며, 출력
회계 증표·이전 증표를 각 트리에 추가하거나 제품 주장을 등록
한다. 같은 트랜잭션에서 감사 기록과 모든 출력의 생성 기록
색인을 기록하고 다음 aid를 증가시킨다. 하나라도 실패하면 전
체 상태 갱신을 취소한다. 따라서 과거 루트의 포함 관계만으로
이미 소비되었거나 현재 동결된 객체를 다시 사용할 수 있도록
정의하지 않는다.

\subsection{아홉 이벤트의 입출력}
Table 3에서 \(N,V,C\)는 각각 회계 증표, 이전 증표, 제품 주장을
뜻한다. \(N_{C}\)는 송신자에게 남는 잔여 회계 증표이다. 각 이벤트
에는 앞 절의 공통 검사가 적용되며, 아래 관계는 그 위에 추가
되는 참조 회계 규칙이다. 등록된 다른 정책을 사용할 수 있더
라도 실제로 검사하지 않는 참조 회계 관계까지 만족한다고 해
석할 수는 없다.

\subsection{초기 등록과 운송}
초기 등록은 입력 없이 ELIGIBLE 회계 증표를 생성하며 \(q>0\),
\(0≤a≤q\), \(e=e_{0}≥0\)을 요구한다. 유효한 초기 탄소 \(e_{0}\)를

% Source PDF page 8, column 1
같고 신뢰된 등록 승인자의 승인이 필요하다.

이전 이벤트는 ELIGIBLE 회계 증표를 소비하여 ELIGIBLE
이전 증표와 같은 역할의 잔여 회계 증표를 생성한다. 두 출력
은 입력 문서 해시를 보존한다. 이전 증표의 송신자와 잔여 회

\begin{equation}
    \begin{aligned}q_{\mathrm{in}}&=q_T+q_C,\quad q_T>0,\quad q_C\ge0,\\a_C&=\left\lfloor\frac{a_{\mathrm{in}}q_C}{q_{\mathrm{in}}}\right\rfloor,\quad a_T=a_{\mathrm{in}}-a_C,\\e_C&=\left\lfloor\frac{e_{\mathrm{in}}q_C}{q_{\mathrm{in}}}\right\rfloor,\\e_T&=e_{\mathrm{in}}-e_C+\delta_T\end{aligned}\tag{13}
\end{equation}
이며 \(δ_{T}≥0\)는 운송 탄소 증분이다. 잔여 회계 증표에는 남는
질량의 기존 탄소 몫을 보존하고, 운송 증분은 보내는 이전 증표
에 더한다. 전량 이전에서도 \(q_{C}=a_{C}=e_{C}=0\)인 잔여 회계
증표를 생성한다.

공개 절대 기한 \(D\)는 이전 증표 안의 값과 같아야 한다. 이전
이벤트 실행 시 검증 계약은 \(D>\) CurrentEpoch\(()\)와 배포 설
정의 허용 범위를 검사한다. 현재 에포크(epoch)는 블록 높이
에 따라 계산하며 증명 생성 시각이나 클라이언트의 로컬 시계
를 사용하지 않는다.

\subsection{인수와 회수}
인수는 수신자의 소유권 비밀을, 회수는 송신자의 소유권 비
밀을 사용하여 이전 증표를 해당 소유자의 회계 증표로 바꾼
다. 두 경로 모두 문서 해시, 역할과 \((q,a,e)\)를 그대로 유지하며
식 (7)의 공통 식별자를 소비한다. 인수에는 기한 조건이 없다.
회수에는 트랜잭션 실행 시

\begin{equation}
    \operatorname{CurrentEpoch}()\le D\tag{14}
\end{equation}
를 요구하므로 기한과 같은 에포크도 포함한다. 기한 전에 만든
증명이라도 기한 이후에 제출하면 거부한다. 기한 이후 미해결
이전 증표는 인수로만 해결할 수 있다. 어느 경로든 현재 소비
여부와 상태 검사를 통과해야 한다.

% Source PDF page 8, column 2

\subsection{병합과 분할}
병합은 서로 다른 두 ELIGIBLE 회계 증표의 소유자와 문서 해
시가 같은 경우에만 허용한다. 출력은 같은 소유자·문서·역할을
가지며

\begin{equation}
    (q,a,e)_{\mathrm{out}}=(q_1+q_2,a_1+a_2,e_1+e_2)\tag{15}
\end{equation}
이다. 같은 회계 증표를 반복 입력하여 합계를 늘리는 것은 중
복 검사에서 거부한다.

분할은 \(q_{in}>0\)인 ELIGIBLE 회계 증표를 둘로 나누며 두 출
력 모두 입력의 소유자·문서·역할을 유지한다. \(q_{in}=q_{1}+q_{2}\)에
대해

\(⌊⌋\)

\begin{equation}
    \begin{aligned}a_2&=\left\lfloor\frac{a_{\mathrm{in}}q_2}{q_{\mathrm{in}}}\right\rfloor,\quad a_1=a_{\mathrm{in}}-a_2,\\e_2&=\left\lfloor\frac{e_{\mathrm{in}}q_2}{q_{\mathrm{in}}}\right\rfloor,\quad e_1=e_{\mathrm{in}}-e_2\end{aligned}\tag{16}
\end{equation}
첫 번째 출력에 잔여량을 주는 순서를 유지한다. 질량 0 출력은
허용하되 그 출력의 크레딧과 탄소도 0이어야 한다.

\subsection{가공}
가공 정책은 정확한 입력·출력 수 \((m,n)\)과 순서 있는 출력 역할
을 고정한다. \(m,n≥1\)이며 모든 입력은 서로 다른 ELIGIBLE
회계 증표이고, 입력과 출력은 모두 공정 수행자 소유다. 지원
하는 유한 입출력 수 집합은 배포 설정에 명시한다. 출력 중 하
나 이상을 ELIGIBLE로 지정하며, 증명자가 수치를 본 뒤 역할
을 바꿀 수 없다.

\(I_{E},I_{W}\) 를 각각 일반용·폐기물 출력의 위치 집합이라 하고,
\(Q_{in}\)을 입력 질량 합, \(Q_{E},Q_{W}\) 를 각 집합의 출력 질량 합이라 하

\begin{equation}
    \begin{gathered}Q_{\mathrm{in}}=Q_E+Q_W+q_{\mathrm{loss}},\\q_{\mathrm{loss}}\ge0,\quad Q_E>0\end{gathered}\tag{17}
\end{equation}

\begin{equation}
    \sum_i a_i^{\mathrm{in}}=\sum_{j\in I_E}a_j^{\mathrm{out}}\tag{18}
\end{equation}
를 만족하며 일반용 출력 사이에서 각 \(0≤a_{j}≤q_{j}\) 범위 안에
비공개로 배분한다. 폐기물 출력의 크레딧은 0이다.

\begin{equation}
    \begin{aligned}E&=\sum_i e_i^{\mathrm{in}}+\delta_P,\quad\delta_P\ge0,\\e_j^{\mathrm{out}}&=\left\lfloor\frac{E q_j^{\mathrm{out}}}{Q_E}\right\rfloor\quad(j\in I_E,\ j\ne j_0),\\e_{j_0}^{\mathrm{out}}&=E-\sum_{j\in I_E\setminus\{j_0\}}e_j^{\mathrm{out}}\end{aligned}\tag{19}
\end{equation}
비례하여 배분하고 반올림 잔여량은 \(j_{0}\)에 준다. 폐기물 출력의
탄소는 0이며, 질량 0인 출력 위치가 있어도 지정된 역할과 잔
여량 수신 순서를 바꾸지 않는다. 이는 폐기물의 물리적 배출량
이 0이라는 뜻이 아니라, 해당 참조 관계에서 탄소를 일반용 출
력에 귀속시키는 규칙이다.

출력 문서 해시는 입력과 달라질 수 있다. 질량 손실과 공정
탄소 증분의 실제 값은 외부 입력이며, 정책이 추가 조건을 검사
할 수 있다. 이 관계 자체가 물질 간 호환성이나 공정의 물리적
수행을 판정하지는 않는다.

% Source PDF page 9, column 1

\subsection{종료와 발급}
종료는 소유자의 회계 증표를 소비하고 후속 객체 없이 종료한
다. ELIGIBLE과 WASTE 모두 허용되며, 폐기물의 유일한 소비
경로이다. 실제 부모는 감사 기록에 암호화하여 남기지만 출력
과 미래 식별자 암호문은 없고 제품 주장도 생성하지 않는다. 모
든 물량의 폐기를 표현할 때도 이 종료 경로를 사용한다.

발급은 \(q>0\)인 ELIGIBLE 회계 증표를 한 번 소비하여 제품
주장을 생성한다. 정책이 정한 재활용 크레딧 하한과 탄소 집약
도 상한을 \(τ_{R}=u_{R}/v_{R}\), \(τ_{C}=u_{C}/v_{C}\)로 두면

\begin{equation}
    v_Ra\ge u_Rq,\quad v_Ce\le u_Cq\tag{20}
\end{equation}
를 검사한다. 분모는 양수이고 배율과 단위는 호환되어야 하며,
비교는 범위가 제한된 정수 교차곱으로 수행한다. 입력 회계 증
표의 \(d\)와 선택한 \(p\), 논스 \(c\)는 식 (9)의 공개 참조값에 결합된다.
회계 증표 소비, 제품 주장 등록, 생성 기록 색인과 실제 부모 기
록은 원자적이다. 같은 회계 증표에서 정책이나 논스를 바꾸어
재발급하거나, 발급 후 종료 또는 종료 후 발급을 수행할 수 없
다.


\begin{table}[htbp]\centering\small
    \setcounter{table}{1}\caption{Verification and authorization routes}
    \begin{tabularx}{\linewidth}{@{}XXX@{}}\toprule
        \textbf{Events} & \textbf{Verifier} & \textbf{Policy and scope}  \\ \midrule
        Ordinary seven  & Fixed per event   & No registry or scope grant \\
        Process         & Registered policy & Policy and scope grant     \\
        Issue           & Registered policy & Policy only                \\
        \bottomrule\end{tabularx}\end{table}

\begin{table}[htbp]\centering\small
    \setcounter{table}{2}\caption{Ordered event inputs and outputs}
    \begin{tabularx}{\linewidth}{@{}XXXX@{}}\toprule
        \textbf{Event} & \textbf{Transition} & \textbf{Parents} & \textbf{Future IDs} \\ \midrule
        Entry          & $\varnothing\to N$  & 0                & 1                   \\
        Transfer       & $N\to(V,N_C)$       & 1                & 2                   \\
        Proceed        & $V\to N_R$          & 1                & 1                   \\
        Recall         & $V\to N_S$          & 1                & 1                   \\
        Merge          & $(N_1,N_2)\to N$    & 2                & 1                   \\
        Split          & $N\to(N_1,N_2)$     & 1                & 2                   \\
        Process        & $N^m\to N^n$        & $m$              & $n$                 \\
        Exit           & $N\to\varnothing$   & 1                & 0                   \\
        Issue          & $N\to C$            & 1                & 0                   \\
        \bottomrule\end{tabularx}\end{table}
